
\section{Matrices}
\subsection{Exercices méthodes :}
\subsubsection{Diagonalisation 1.}
\subsubsection{Diagonalisation 2.}
\subsubsection{Diagonalisation 3.}

\subsubsection{Utiliser la Trace 1.}
\subsubsection{Utiliser la Trace 2.}
\subsubsection{Utiliser la Trace 3.}

\subsection{Exercices classiques.}
Soit $A,B \in \mathcal{M}_n(\mathbb{R})$, et $\lambda_0,\dots,\lambda_n \in \mathbb{K}$
telles 
\subsubsection{Nilpotence implique trace nulle.}
\subsubsection{Lemme de Hadamard.}
\subsubsection{Conditions pour que $rg(f)\leq rg(g)$, (\textit{X})}
\textbf{\underline{Énoncé :}}
Soit $E,F$ deux $\mathbb{K}$-espaces vectoriels non nuls de dimension finie, 
$f$ et $g$ dans $\mathcal{L}(E,F)$. 

Montrer que $rg(g)\leq rg(f)$ si et seulement s'il existe 
$h \in GL(F)$ et $k \in \mathcal{L}(E)$ tels que $h \circ g = f \circ k$

\subsubsection{Fonctions multiplicatives et inversibilité, (\textit{X})}
\textbf{\underline{Énoncé :}} 
Soit $f \colon \mathcal{M}_n(\mathbb{K})\rightarrow \mathbb{K}$ non constante 
telle que : $\forall A,B \in \mathcal{M}_n(\mathbb{K}), f(AB)=f(A)f(B)$.

Montrer que $f(A)=0$ si et seulement si $A$ n'est pas inversible.
\newpage

\subsubsection{Décomposition d'une matrice par le rang}
\textbf{\underline{Énoncé :}} 
Soit $A \in \mathcal{M}_n(\mathbb{K})$ de rang $r$, 

Montrer qu'il existe $r$ matrices de rang $1$ dont 
la somme fait $A$.      

\textbf{\underline{Solution :}}

Sans perte de généralité, on suppose que les $r$ premières colonnes de $A$ sont libres. 

C'est possible car il y a $r$ colonnes libres comme $A$ est de rang $r$. 
et on simplifie seulement nos indices. 

Alors : $$A = 
\begin{pmatrix} 
 &C_1 &C_2 &\dots &C_r & C_{r+1} &\dots&C_n \\ 
\end{pmatrix}
$$
Il se trouve que pour tout $i> r$, $C_i \in \text{Vect}(C_1,\dots C_r)$.

Donc pour tout $i > r$, il existe $a_{i,1},\dots,a_{i,r}$ tels que 
$\displaystyle C_i = \sum_{k=1}^{r}a_{i,k}C_k $
On pose pour tout $i \in \llbracket 1,r \rrbracket$

$$ M_i = \begin{pmatrix} &0&\dots&\underbrace{C_i}_{\text{i-éme colonne}} &\dots &\underbrace{a_{j,i}}_{\text{j-éme colonne}} &\dots &a_{n,i}\end{pmatrix}$$


\subsubsection{Problème 12 :}
\textbf{\underline{Énoncé :}} 

Soit $A_1,\dots, A_r$ des parties de $\llbracket 1,n\rrbracket$.

On suppose qu'il existe $b<a$ tels que pour tout $i,j \in \llbracket 1,r\rrbracket$ : 

$$ 
\text{Card}(A_i \bigcap A_j) = 
\begin{cases} a & \mbox{si }  i=j \\ b & \mbox{si } i \neq j \end{cases}
$$

Montrer que $r \leq n$.

\subsubsection{Décomposition en inverse}
\textbf{\underline{Énoncé :}} 

Soit $\mathbb{K}$ un corps de caractéristique $\neq 2$

Soit $M \in \mathcal{M}_n(\mathbb{K})$, 
montrer que $M$ s'écrit comme la somme de deux matrices inversibles.

\textbf{\underline{Solution :}}

Les coefficients diagonaux de $M$ sont $d_1,\dots,d_n$.

Alors pour tout $i \in \llbracket 1,n\rrbracket$, on pose 

$$
a_i = \begin{cases} \frac{d_i}{2} & \mbox{si }  d_i\neq 0 \\ 1 & \mbox{sinon } \end{cases} \text{ et } b_i = \begin{cases} \frac{d_i}{2} & \mbox{si }  d_i\neq 0 \\ -1 & \mbox{sinon }  \end{cases}
$$

On a alors 
$$
M =
\begin{pmatrix}
d_i &  &(S)\\
    &\ddots & \\
 (I)   &  & d_n 
\end{pmatrix}
$$

On pose 
$$
A = \begin{pmatrix}
a_i &  &(0)\\
    &\ddots & \\
 (I)   &  & a_n 
\end{pmatrix} 
\text{ et } 
B = 
\begin{pmatrix}
b_i &  &(S)\\
    &\ddots & \\
 (0)   &  & b_n 
\end{pmatrix}
$$

Par définition, on a bien $A + B = M$

Et $A$ et $B$ sont bien inversibles car triangulaires à coefficients diagonaux non nuls. $\square $

